% Versão de trabalho: 28/09/2026. Documento principal: TCCFinal.tex.
% Compilação: pdfLaTeX + BibTeX. Métodos propostos; experimentos pendentes.
\documentclass{UNESC}
\usepackage[T1]{fontenc}
\usepackage[utf8]{inputenc}
\renewcommand{\rmdefault}{phv}
\renewcommand{\sfdefault}{phv}
\usepackage[brazilian]{babel}
\usepackage{graphicx}
\usepackage{amsmath}
\usepackage{array}
\usepackage{booktabs}
\usepackage{tabularx}
\usepackage{enumitem}
\usepackage{float}
\usepackage[hang,flushmargin]{footmisc}
\usepackage{setspace}
\usepackage[labelsep=space]{caption}
\DeclareCaptionLabelFormat{comtraco}{#1~#2 -}
\captionsetup{labelformat=comtraco}
\usepackage{hyperref}
\usepackage[alf,recuo=0.0cm,abnt-emphasize=bf]{abntex2cite}
\hypersetup{setpagesize=false,colorlinks=true,urlcolor=blue,linkcolor=black,citecolor=black}
\urlstyle{same}
\setlength{\emergencystretch}{2em}

\title{MONITORAMENTO DE PERDAS POR SUJIDADE EM SISTEMAS FOTOVOLTAICOS COM APRENDIZADO DE MÁQUINA}
\author{Matheus Luís Mantovani Passoni\footnote{Curso de Ciência da Computação, Universidade do Extremo Sul Catarinense (Unesc), Criciúma - Santa Catarina - Brasil. matheus.mantovani28@gmail.com},
Allan Farias Fávaro\footnote{Orientador, Curso de Ciência da Computação, Universidade do Extremo Sul Catarinense (Unesc), Criciúma - Santa Catarina - Brasil.}}
% Acrescentar o e-mail do orientador após confirmação.

\begin{document}
\onehalfspacing
\maketitle
\noindent\textit{Versão de trabalho: proposta metodológica em desenvolvimento, sem resultados experimentais.}
\vspace{0.7em}

\noindent\textbf{Resumo:} O acúmulo de sujidade sobre módulos fotovoltaicos pode reduzir sua geração, mas a identificação dessas perdas exige considerar as condições ambientais e outras alterações operacionais. Este trabalho propõe desenvolver e avaliar um método de monitoramento baseado em aprendizado de máquina para estimar a potência esperada em uma condição de referência de baixa sujidade e analisar desvios persistentes em relação à potência medida. A pesquisa terá abordagem quantitativa, com desenvolvimento de um artefato computacional e experimentos comparativos. Serão investigados registros públicos de irradiância, temperatura e potência, com seleção documentada dos períodos de referência. Um modelo Random Forest será comparado com uma regressão de referência, mediante separação temporal dos dados. A avaliação contemplará erros de estimativa de potência e, quando houver rótulos independentes ou cenários controlados, a qualidade dos alertas. Experimentos semissintéticos poderão complementar a análise, com identificação explícita das perdas introduzidas artificialmente. Nesta versão, apresentam-se a fundamentação e o protocolo proposto; resultados, conclusões e validação em campo permanecem pendentes.

\vspace{0.5em}
\noindent\textbf{Palavras-chave:} energia fotovoltaica; sujidade; aprendizado de máquina; regressão; monitoramento de desempenho.

\vspace{0.8em}
\noindent\textbf{Abstract:} Soiling on photovoltaic modules can reduce power generation, but identifying these losses requires accounting for environmental conditions and other operational changes. This study proposes the development and evaluation of a machine learning monitoring method to estimate expected power under a low-soiling reference condition and analyze persistent deviations from measured power. The research will follow a quantitative approach, combining a computational artifact with comparative experiments. Public irradiance, temperature, and power records will be investigated, with documented selection of reference periods. A Random Forest model will be compared with a regression baseline using chronological data partitions. Evaluation will address power estimation errors and, whenever independent labels or controlled scenarios are available, alert performance. Semi-synthetic experiments may complement the analysis, with explicit identification of artificially introduced losses. This working version presents the rationale and proposed protocol; experimental results, conclusions, and field validation remain pending.

\vspace{0.5em}
\noindent\textbf{Keywords:} photovoltaic energy; soiling; machine learning; regression; performance monitoring.

\section{INTRODUÇÃO}

A geração fotovoltaica depende da radiação incidente e das condições de operação dos módulos, incluindo sua temperatura \cite{king2004}. Entre os fatores que podem reduzir seu desempenho encontra-se a sujidade, denominada \textit{soiling}, caracterizada pela deposição de poeira e outros contaminantes sobre sua superfície. Seu monitoramento pode apoiar decisões de manutenção, mas envolve a escolha de medições e procedimentos adequados para estimar as perdas \cite{bessa2021}.

O problema não se resume a observar se houve redução da geração. A modelagem de potência permite relacionar as condições ambientais à saída elétrica esperada e comparar essa estimativa com as medições para acompanhar a operação \cite{li2024}. Para o presente estudo, será necessário estabelecer uma referência compatível com a irradiância e a temperatura observadas, de modo que o desvio não represente apenas uma diferença entre condições meteorológicas.

Potência e energia constituem grandezas diferentes: o watt (W) equivale a um joule por segundo, enquanto o joule (J) é unidade de energia \cite{bipm2026}. Neste trabalho, a potência será expressa em W ou kW e sua integração temporal em Wh ou kWh, sendo $1\,\mathrm{Wh}=3600\,\mathrm{J}$. O indicador proposto representará desempenho relativo à referência estimada; não será uma medição direta da eficiência física de conversão de um módulo.

Na revisão de \citeonline{bessa2021}, as estratégias de monitoramento incluem estações de sujidade, sensores ópticos e análise de imagens, além de abordagens de modelagem. Essa diversidade mostra que a escolha de uma solução depende tanto da informação necessária quanto dos recursos disponíveis. No recorte deste TCC, o interesse estará nos registros numéricos de operação e ambiente, sem exigir processamento de imagens ou a construção de um equipamento de medição de sujeira.

Uma referência física de potência pode ser construída com parâmetros do sistema. O manual de \citeonline{dobos2014} descreve os submodelos do PVWatts e suas hipóteses. Na implementação documentada pelo pvlib, a potência em corrente contínua combina irradiância efetiva, potência nominal e correção térmica \cite{pvlib_pvwatts}. Será necessário compatibilizar a saída com o ponto de medição: o PVWatts trata separadamente a geração em corrente contínua e sua conversão pelo inversor em corrente alternada \cite{dobos2014}.

Outra possibilidade é estimar essa relação a partir de observações por regressão supervisionada, com variáveis de entrada e um alvo numérico \cite{pedregosa2011}. Neste trabalho, as entradas principais serão irradiância e temperatura, enquanto o alvo será a potência dos períodos de referência. A expressão aprendizado de máquina designará esse ajuste computacional. Como antecedente aplicado, \citeonline{li2024} comparam modelos de aprendizado de máquina, incluindo Random Forest, com alternativas físicas para estimar a potência de sistemas degradados.

Entre os algoritmos disponíveis, o Random Forest reúne árvores de decisão e agrega suas estimativas. Sua formulação admite problemas de regressão \cite{breiman2001}. Na implementação do scikit-learn, a previsão resulta da combinação das saídas de árvores treinadas com amostras dos dados \cite{sklearn_rf}. A escolha inicial para este estudo decorre da possibilidade de investigar relações não lineares entre variáveis tabulares com uma implementação acessível. Isso constitui uma decisão metodológica a avaliar, sem pressupor que o algoritmo será superior às alternativas.

A estimativa de sujidade por dados operacionais já possui antecedentes científicos. \citeonline{deceglie2018} apresentam o método \textit{stochastic rate and recovery} (SRR), que identifica intervalos de sujidade e constrói perfis estocásticos a partir do desempenho observado, sem exigir dados de precipitação. O estudo inclui validação com medições de estações de sujidade. Essa referência reforça a necessidade de comparar a estimativa do algoritmo com evidência externa quando se pretende atribuir uma perda à sujeira.

\citeonline{skomedal2020} abordam a coexistência de perdas, propondo decompor um índice de desempenho em componentes de sujidade, degradação e sazonalidade. Já \citeonline{mejia2013} analisam sistemas residenciais e comerciais e relacionam a recuperação após chuva à quantidade precipitada. Esses trabalhos motivam dois cuidados deste projeto: não atribuir qualquer desvio à sujidade e não assumir que todo evento de chuva restabelece uma condição limpa.

No contexto brasileiro, \citeonline{Tramontin_Rampinelli_2024} analisam dois sistemas no Rio Grande do Sul, utilizando dados de inversores antes e depois da limpeza e informações de estações meteorológicas. Os autores empregam regressão linear na relação entre irradiação e energia e discutem limitações dos dados disponíveis. Para este TCC, o estudo contribui ao relacionar análise computacional e intervenção documentada, sem justificar a transferência de seus ganhos percentuais para outra instalação.

No campo da regressão, \citeonline{kalimeris2023data} investigam a estimativa da razão de sujidade com dados de potência, irradiância e temperatura. Seus procedimentos incluem reconhecer eventos potenciais de limpeza e selecionar períodos de aprendizagem. O trabalho utiliza regressão Ridge com atributos polinomiais; em uma das abordagens, a informação de chuva permite operar sem registros obrigatórios de limpeza manual. Assim, não seria adequado caracterizar a literatura como restrita a fórmulas estáticas ou afirmar que o uso de aprendizado de máquina constitui, por si só, uma novidade.

Uma pesquisa mais recente, de \citeonline{tsanakas2025}, combina duas perspectivas: análise retrospectiva das perdas a partir do desempenho elétrico e previsão de perdas com variáveis ambientais. A proposta evidencia a distinção entre estimar o estado de operação e antecipar sua evolução. No presente TCC, o primeiro objetivo será estimar a potência de referência no instante das medições e acompanhar desvios persistentes, mantendo o escopo compatível com a base que puder ser obtida.

Esse recorte também delimita o uso do termo previsão. \citeonline{li2024} distinguem o monitoramento com dados meteorológicos medidos da previsão de potência futura com dados meteorológicos previstos. Neste projeto, o sistema estimará a potência de referência no instante observado. Antecipar a geração das próximas horas exigiria outro protocolo de entradas e horizonte, não incluído no experimento inicial.

A análise desses antecedentes orienta a contribuição pretendida: avaliar uma implementação de complexidade controlada, que compare modelos de referência e examine a repercussão de seus erros na qualidade dos alertas. A questão a investigar no recorte escolhido é se a maior flexibilidade do Random Forest produz ganho útil para o monitoramento, quando comparado com uma regressão simples nas mesmas condições. Essa questão constitui uma delimitação experimental; uma afirmação de ineditismo geral dependeria de revisão mais abrangente da literatura.

A preparação da base terá papel central nessa avaliação. A seleção de períodos de aprendizagem próximos a eventos de limpeza faz parte da abordagem de \citeonline{kalimeris2023data}. Neste projeto, os períodos de referência deverão ser justificados por evidências e representar as condições ambientais de interesse. Sem essa justificativa, a interpretação ficará limitada a mudanças em relação ao histórico operacional, sem comprovar a potência de um sistema perfeitamente limpo.

A interpretação da diferença entre potência estimada e medida deverá considerar a coexistência de perdas discutida por \citeonline{skomedal2020} e a incerteza da modelagem. Por isso, propõe-se avaliar uma regra de persistência temporal e cenários que possam ser confundidos com sujidade. A regra será uma escolha experimental deste trabalho, sem pressupor que a persistência identifique a causa do desvio.

Outro aspecto é a origem dos dados. O repositório PVDAQ disponibiliza registros de desempenho e metadados de diferentes instalações, com sensores ambientais em parte dos conjuntos \cite{pvdaq2021}. Ele será uma fonte candidata para a pesquisa. A seleção não poderá se apoiar apenas no volume de registros: a correspondência temporal entre potência e ambiente, a disponibilidade de unidades e a qualidade da referência deverão ser verificadas.

Quando não houver rótulos de sujidade suficientes, cenários controlados poderão complementar a avaliação. O exemplo do RdTools para o sistema 4 do PVDAQ utiliza um sinal de sujidade sintetizado, explicitamente identificado como tal em sua documentação \cite{rdtools_exemplo}. No presente trabalho, uma combinação de medições com perdas introduzidas artificialmente será denominada semissintética. Seus resultados permitirão avaliar perdas conhecidas naquele experimento, sem substituir uma validação de sujidade em campo.

Também será preservada a ordem temporal dos registros na avaliação. \citeonline{cerqueira2020} mostram que a estimativa de desempenho depende do protocolo e das características da série, com resultados favoráveis à avaliação fora da amostra nos cenários não estacionários analisados. Para representar o uso pretendido neste TCC, o treinamento antecederá o teste. O controle abrangerá ainda o vazamento de dados: informação indevida nas entradas ou na preparação pode produzir avaliações excessivamente otimistas \cite{kapoor2023}.

Diante desse contexto, formula-se a questão de pesquisa: em que medida um modelo Random Forest melhora a estimativa da potência de referência e apoia a identificação de perdas persistentes compatíveis com sujidade, em comparação com uma regressão de referência? A hipótese é que o modelo reduza o erro de estimativa e melhore a detecção nos cenários avaliados, sem aumento relevante dos falsos alertas. Essa hipótese será testada e poderá ser rejeitada.

O objetivo geral é desenvolver e avaliar um método computacional de monitoramento de sistemas fotovoltaicos baseado na estimativa de potência por aprendizado de máquina e na análise de desvios persistentes associados à sujidade. Os objetivos específicos são:
\begin{enumerate}[label=\alph*),leftmargin=1.2cm,itemsep=0.1em]
\item selecionar e preparar dados temporais de potência e ambiente, documentando sua origem, qualidade e condição de referência;
\item implementar e comparar Random Forest e regressão Ridge na estimativa da potência esperada;
\item desenvolver um indicador de desvio e uma regra de alerta com persistência temporal;
\item avaliar separadamente a estimativa de potência e a detecção de perdas, distinguindo observações reais e cenários simulados;
\item implementar um protótipo que apresente as séries, os indicadores e os alertas.
\end{enumerate}

A justificativa está na necessidade de compreender quando uma estimativa aprendida a partir de dados oferece informação útil ao acompanhamento fotovoltaico. A contribuição computacional pretendida será uma avaliação reprodutível dos modelos e dos alertas, com análise de limitações, além da construção do protótipo. A aplicação à microgeração é uma motivação prática, condicionada à disponibilidade de sensores e à representatividade da base. O método não pretende, nesta etapa, determinar automaticamente a causa de qualquer perda nem o momento economicamente ótimo para limpeza.

\section{MATERIAIS E MÉTODOS}

\subsection{DELINEAMENTO E DADOS}

A pesquisa será aplicada e quantitativa, com desenvolvimento e avaliação experimental de um artefato computacional. A organização do trabalho se apoia na relação entre construção e avaliação de artefatos discutida por \citeonline{hevner2004}, sem reivindicar nesta etapa a execução completa de um método de \textit{design science}. Os resultados próprios decorrerão da comparação dos métodos com dados secundários. Por delimitação experimental, será analisada uma instalação de cada vez.

A base definitiva ainda não foi selecionada. A primeira fonte candidata é o PVDAQ \cite{pvdaq2021}. A alternativa controlada é o conjunto do sistema 4 utilizado pelo RdTools, cuja documentação apresenta potência, irradiância no plano dos módulos, temperatura ambiente e vento, além de um sinal de sujidade sintetizado \cite{rdtools_exemplo}. A disponibilidade efetiva e a qualidade de cada variável deverão ser verificadas no arquivo adquirido.

Após a seleção, serão registrados identificação do sistema, endereço e licença da fonte, data de obtenção, período, intervalo amostral, unidades e quantidade de registros antes e depois dos filtros. Cada variável será classificada como medida, estimada ou simulada. A referência e o alvo deverão representar o mesmo ponto de medição, em corrente alternada ou contínua.

\begin{table}[htbp]
\centering
\caption{Variáveis previstas e sua função}
\label{tab:variaveis}
\small
\begin{tabularx}{\linewidth}{@{}p{3.5cm}X@{}}
\toprule
Variável & Uso no experimento \\
\midrule
Data e hora & Ordenação, sincronização e divisão temporal. \\
Irradiância no plano dos módulos & Entrada principal, em W/m$^2$. \\
Temperatura & Entrada; distinguir módulo, célula e ambiente. \\
Potência medida & Alvo no treinamento de referência e comparação no monitoramento, em W ou kW. \\
Vento & Entrada opcional, conforme os dados disponíveis. \\
Chuva e limpeza & Evidência contextual para interpretar mudanças e avaliar a referência. \\
Perda simulada & Controle e avaliação; excluída das entradas do modelo. \\
\bottomrule
\end{tabularx}
\par\smallskip\raggedright\footnotesize Fonte: elaboração própria, com base em \citeonline{king2004}, \citeonline{rdtools_exemplo} e \citeonline{kapoor2023}.
\end{table}

O modelo de \citeonline{king2004} distingue temperatura ambiente, do módulo e da célula, além de tratar a irradiância incidente no plano do módulo. Por isso, temperatura ambiente e irradiância horizontal não serão substituídas diretamente por grandezas diferentes. Se necessário, será usada uma transformação física documentada. Como prevenção de vazamento por variáveis que reproduzem o alvo \cite{kapoor2023}, a potência contemporânea e suas reproduções diretas, como tensão multiplicada por corrente, serão excluídas das entradas do estimador.

\subsection{PREPARAÇÃO E CONDIÇÃO DE REFERÊNCIA}

O tratamento de horários e valores atípicos integra a preparação dos dados em \citeonline{li2024}; o fluxo do RdTools também documenta filtros de qualidade para a análise de desempenho \cite{rdtools_exemplo}. Neste projeto, serão tratados duplicidades, fuso, unidades e ausências, com identificação de noite, baixa irradiância, indisponibilidade, erros físicos e limitação do inversor. Os critérios e as exclusões serão registrados. Por escolha do protocolo, potência baixa não será descartada apenas por divergir da previsão.

Serão priorizados períodos sustentados por limpeza documentada, medição independente ou informação técnica sobre a base, considerando os procedimentos de \citeonline{deceglie2018} e \citeonline{kalimeris2023data}. \citeonline{mejia2013} adotam um critério de precipitação específico às condições estudadas; ele não será transferido automaticamente para outra instalação. Sem referência suficiente, as conclusões se limitarão a perdas adicionais. A qualidade dos sensores \cite{bessa2021} e a cobertura das condições ambientais integrarão a análise de validade.

Como adaptação experimental da estratégia de sinal sintetizado apresentada pelo RdTools \cite{rdtools_exemplo}, a série original será a referência operacional, sem ser descrita como perfeitamente limpa. Perdas adicionais serão introduzidas na validação e no teste. O erro de regressão será calculado contra a série original reservada, e os alertas sobre a série modificada. Para manter a separação entre entradas e informação de avaliação \cite{kapoor2023}, o detector não receberá os fatores artificiais nem a potência original reservada do teste.

\subsection{MODELAGEM E VALIDAÇÃO TEMPORAL}

Adotando o mapeamento entre variáveis ambientais e potência usado na modelagem fotovoltaica \cite{li2024}, define-se, para este experimento,
\begin{equation}
\widehat{P}_{\mathrm{ref}}(t)=f\bigl(G(t),T(t),\mathbf{z}(t)\bigr),
\end{equation}
em que $G$ representa irradiância, $T$ a temperatura escolhida e $\mathbf{z}$ eventuais entradas adicionais documentadas. O alvo de treinamento será a potência nos períodos de referência. Novas entradas dependerão de disponibilidade no instante de uso e avaliação na validação.

O experimento comparará Random Forest \cite{breiman2001} com regressão Ridge. Esta última introduz regularização na estimação dos coeficientes, conforme a formulação de \citeonline{hoerl1970}. Serão usados irradiância, temperatura e sua interação, com padronização ajustada no treinamento. A regularização e os hiperparâmetros da floresta serão escolhidos nos dados de desenvolvimento. Para a floresta, serão investigados número de árvores, profundidade e mínimo de observações nas folhas, parâmetros documentados na API \cite{sklearn_rf}. O PVWatts poderá complementar a comparação quando houver metadados compatíveis \cite{dobos2014}.

O protocolo utilizará blocos cronológicos de treinamento, validação e teste, apoiado na discussão de \citeonline{cerqueira2020}. O período mais recente será reservado à avaliação final. Nos períodos anteriores serão ajustados modelos e alertas, com janelas temporais. A documentação de TimeSeriesSplit exige amostras igualmente espaçadas para comparar durações dos testes \cite{sklearn_temporal}; caso os filtros gerem irregularidade, serão definidos blocos por datas, com duração comparável. As datas dependerão da cobertura da base.

Conforme a orientação técnica para prevenir vazamento, imputações, padronizações e demais transformações serão ajustadas somente no treinamento de cada divisão \cite{sklearn_praticas}. O protocolo excluirá interpolação com informação futura, registrará sementes e manterá o teste fora da seleção dos candidatos, em consonância com os riscos de avaliação discutidos por \citeonline{kapoor2023}.

\subsection{INDICADOR E ALERTAS}

A comparação entre potência observada e uma referência modelada é empregada na análise de sujidade por \citeonline{kalimeris2023data}. Neste projeto, será adotado o complemento da razão entre essas potências, com a notação
\begin{equation}
d(t)=1-\frac{P_{\mathrm{obs}}(t)}{\widehat{P}_{\mathrm{ref}}(t)}.
\end{equation}
O cálculo exigirá previsão positiva e irradiância adequada; a instabilidade de medidas relativas com denominadores próximos de zero é discutida por \citeonline{hyndman2006}. Por definição da equação, valores negativos indicam potência observada superior à referência e serão preservados. Considerando a coexistência de perdas \cite{skomedal2020}, o indicador não será denominado perda de sujidade medida sem evidência adicional.

Propõe-se agregar essa razão diariamente, ponderando pela energia de referência nos mesmos intervalos válidos. Trata-se de uma definição operacional deste estudo:
\begin{equation}
D_j=1-\frac{\sum_{t\in V_j}P_{\mathrm{obs}}(t)\Delta t}
{\sum_{t\in V_j}\widehat{P}_{\mathrm{ref}}(t)\Delta t},
\end{equation}
em que $V_j$ contém os intervalos válidos do dia $j$ e $\Delta t$ é sua duração, com potência média no intervalo. A soma aproxima a integração da potência; com kW e horas, resulta em kWh, de acordo com a relação entre energia e potência \cite{bipm2026}. O protocolo exigirá cobertura mínima diária e não extrapolará registros incompletos como geração integral do dia.

Como regra proposta neste estudo, o alerta dependerá de o indicador superar um limiar por dias consecutivos válidos. Limiar, persistência e eventual média móvel calculada apenas com dados presentes e passados serão escolhidos na validação e fixados antes do teste, preservando sua independência \cite{kapoor2023}. Lacunas interromperão a contagem. A saída apresentará intensidade, duração e qualidade dos dados, orientando inspeção. Não se atribui à literatura um limiar universal nem eficácia já demonstrada dessa regra.

\subsection{EXPERIMENTOS E MÉTRICAS}

A estimativa de potência será avaliada por erro absoluto médio (MAE), raiz do erro quadrático médio (RMSE) e coeficiente de determinação ($R^2$), utilizando as definições documentadas pelo scikit-learn \cite{sklearn_metricas}. Para $n$ pares de valores observados $y_i$ e estimados $\widehat y_i$, com média observada $\overline y$, serão calculados
\begin{equation}
\mathrm{MAE}=\frac{1}{n}\sum_{i=1}^{n}|y_i-\widehat y_i|,\qquad
\mathrm{RMSE}=\sqrt{\frac{1}{n}\sum_{i=1}^{n}(y_i-\widehat y_i)^2},
\end{equation}
\begin{equation}
R^2=1-\frac{\sum_i(y_i-\widehat y_i)^2}{\sum_i(y_i-\overline y)^2}.
\end{equation}
MAE e RMSE terão a unidade da potência; $R^2$ será adimensional, poderá ser negativo e não será interpretado em subconjuntos com alvo constante \cite{sklearn_metricas}. Os modelos serão comparados nos mesmos registros, também por período e faixa de irradiância. O erro percentual absoluto médio não será a medida principal, devido aos problemas de divisão por valores nulos ou próximos de zero discutidos por \citeonline{hyndman2006}.

Quando houver rótulos independentes ou cenários controlados, serão calculados precisão, sensibilidade e F1 \cite{sklearn_metricas}. Sendo VP, FP e FN os números de verdadeiros positivos, falsos positivos e falsos negativos,
\begin{equation}
\begin{aligned}
\text{Precisão}&=\frac{VP}{VP+FP},\qquad
\text{Sensibilidade}=\frac{VP}{VP+FN},\\
F1&=\frac{2VP}{2VP+FP+FN}.
\end{aligned}
\end{equation}
Denominadores nulos serão sinalizados, junto às contagens. A escolha de precisão e sensibilidade considera a discussão de \citeonline{saito2015} sobre classes desbalanceadas; sua aplicação aqui será avaliada conforme a frequência real dos eventos. Por definição deste protocolo, também serão apresentados falsos alertas por dia válido e atraso entre início do evento e primeiro alerta. A unidade de avaliação e a rotulagem serão fixadas antes do teste. Sem rótulos independentes, a análise de sujidade real será descritiva.

Nos cenários semissintéticos, será adotada uma transformação multiplicativa, inspirada no uso de um fator artificial de sujidade no exemplo do RdTools \cite{rdtools_exemplo}:
\begin{equation}
P_{\mathrm{sim}}(t)=P_{\mathrm{base}}(t)\,s(t),\qquad 0\leq s(t)\leq1,
\end{equation}
mantendo a série original e a transformação aplicada. Por escolha experimental, serão considerados controle sem perda adicional, perdas graduais e recuperação simulada. Cenários de indisponibilidade, falha de sensor e redução abrupta examinarão confusões. O fator imposto representa apenas uma redução proporcional de potência; não modela a distribuição espacial dos depósitos nem todos os efeitos elétricos discutidos na literatura sobre sujidade \cite{bessa2021}. Os cenários de teste serão separados dos usados na calibração.

A comparação empregará as mesmas entradas e partições, relacionando erro de previsão, falsos alertas e atraso de detecção. Resultados semelhantes ou superiores da regressão simples também serão reportados. Evidências de simulação não serão apresentadas como comprovação de eficácia de limpeza em uma instalação real.

\subsection{PROTÓTIPO E REPRODUTIBILIDADE}

A implementação está prevista em Python, com pandas e NumPy para preparação e scikit-learn para modelagem \cite{pedregosa2011}. O protótipo inicial receberá arquivos CSV e apresentará potência medida e prevista, desvio diário, cobertura e alertas. Uma interface local em Streamlit poderá ser acrescentada após a validação do processamento.

Serão preservados dados brutos, transformações, configurações, versões e sementes, documentando as etapas para permitir a conferência dos riscos metodológicos discutidos por \citeonline{kapoor2023}. Como decisão do experimento, treinamento e monitoramento terão etapas separadas e o modelo permanecerá fixo durante o teste. A utilização em tempo real e em outras instalações dependerá de nova avaliação.

\section{RESULTADOS E DISCUSSÃO}
Os experimentos ainda não foram executados. Esta seção será preenchida com a caracterização da base, a comparação dos modelos, a avaliação dos alertas e a discussão das limitações e dos trabalhos relacionados. Nenhuma medida de desempenho ou economia é apresentada nesta versão.
% Inserir: qualidade da base; MAE/RMSE/R2; curvas do teste; falsos alertas;
% atraso por evento; confusões; relação crítica com pesquisas anteriores.

\section{CONCLUSÃO}
As conclusões dependerão dos resultados e deverão responder à questão de pesquisa, retomando os objetivos e indicando se a hipótese foi sustentada. Serão explicitadas contribuições, limitações, riscos à validade e possibilidades de novas investigações. A aplicação à microgeração e a identificação de sujidade real permanecerão condicionadas à representatividade da base e à existência de evidência independente.

\section*{DECLARAÇÃO DE USO DE INTELIGÊNCIA ARTIFICIAL}
Ferramentas de inteligência artificial generativa foram utilizadas como apoio à organização textual, ao levantamento preliminar de fontes e à discussão de alternativas metodológicas. A conferência das referências, as decisões de pesquisa e a responsabilidade pelo conteúdo cabem aos autores. Não foram produzidos resultados experimentais para esta versão.

\bibliographystyle{abntex2-unesc}
\renewcommand{\refname}{REFERÊNCIAS}
\bibliography{referencias_final}
\end{document}
